% !PRINTED 317
\weberpage{317}

To see this, we denote, in the earlier notation, 1, 2, 3 by $(\xi)$, $(\eta)$, $(\zeta)$; since these three quantities form a triple, there is [by \S.\,81, (5)] a system $(\alpha)$ such that
\weq{4}{\Sigma\alpha\xi \equiv 0, \quad \Sigma\alpha\eta \equiv 0, \quad \Sigma\alpha\zeta \equiv 0 \ (\mathrm{mod}\ 2)}
holds. If we now apply a substitution from $Q$:
\[
\begin{pmatrix} \xi_1, & \xi_2, & \xi_3 \\ \xi_1', & \xi_2', & \xi_3' \end{pmatrix},
\]
in which we may take $\xi_i' = \xi_i + h_i$, then it follows from (4):
\[
\Sigma\alpha\xi' \equiv \Sigma\alpha\eta' \equiv \Sigma\alpha\zeta' \equiv \Sigma\alpha h,
\]
that is, the $(\xi')$, $(\eta')$, $(\zeta')$ form the triple belonging to $(\alpha)$ if $(h)$ itself belongs to this triple, or, in the other case, $(h)$, $(\xi')$, $(\eta')$, $(\zeta')$, is the quadruple belonging to $(\alpha)$.

If now again $k$ denotes any substitution from $P$ and $\gamma$ a substitution from $Q$ by which the quadruple 0, 1, 2, 3 passes into 4, 5, 6, 7, then we can determine the elements $\tau'$, $\tau''$ from $T_{24}$ and $\gamma'$, $\gamma''$ from $Q$ so that
\[
k = \tau'\gamma' = \gamma^{-1}\tau''\gamma'',
\]
whereby, as above, it is proved that the two systems (3) are imprimitively connected by $k$.

After adjunction of a square root to the domain of rationality, the quantities 0, 1, 2, 3 are then the roots of a biquadratic equation.

By these considerations, for primitive equations of degree 8 the possibility is also excluded that one of the seven quantities $\Psi_{\xi_1,\xi_2,\xi_3}$ [\S.\,80, (1)] vanishes. For these functions are rational functions of the roots $X$; hence if one of them vanishes, every permutation of the group of the equation can be applied to the rational equation $\Psi = 0$. From this it follows, by the theorem \S.\,80, 2., that either all seven functions $\Psi$ must vanish, in which case the roots would be rational, or that the substitutions $\varrho^{-1}$, applied to the $(\xi)$, must form an intransitive group. The degree of the group could then not be divisible by 7. Now the group of the $\varrho^{-1}$ is isomorphic with the group $T$ of the $\sigma$ (\S.\,37, 8.), and therefore the degree of $T$ is also not divisible by 7, and consequently
%CONT>>
